\documentclass[11pt]{article}

\usepackage[a4paper,margin=2.5cm]{geometry}
\usepackage[T1]{fontenc}
\usepackage[utf8]{inputenc}
\usepackage{lmodern}
\usepackage{hyperref}
\usepackage{enumitem}
\usepackage{verbatim}
\usepackage{framed}
\usepackage{tabularx}
\usepackage{url}

% Add spacing between paragraphs
\setlength{\parskip}{0.5em}
%\setlength{\parindent}{0pt}

\title{SÁL242F Exercise 4: Simulation of Questionable Research Practices}
\author{}
\date{}

\begin{document}
\maketitle

\section*{Overview}

In this exercise you will use the simulation app at \url{https://sim.likan.is} to explore how statistical results can be influenced by randomness, researcher choices, and questionable research practices. The main goal is to build intuition by experimenting with the simulations and observing what happens under different settings.

\section*{Instructions}

Follow the instructions in the web app step by step. Explore the simulations extensively and try many different settings so that you begin to develop an intuition for the phenomena being illustrated.

\noindent
Your task is to:
\begin{itemize}
    \item work through the app step by step,
    \item vary the settings and rerun simulations many times,
    \item pay attention to when results look convincing and when they seem to arise mainly from chance or flexibility in analysis,
    \item note patterns, surprises, and anything that seems misleading or unrealistic,
    \item prepare to discuss your observations in class.
\end{itemize}

\noindent
In class, we will discuss the following questions:
\begin{itemize}
    \item What did you observe about p-values, significance, and repeated sampling?
    \item What kinds of researcher choices seemed to have the biggest effect on the results?
    \item What helped you build intuition about p-hacking or other questionable research practices?
    \item What are the limitations of these simulations?
    \item How far can we extrapolate from the findings in the app to real research practice?
    \item What changes, extensions, or new simulations would you like to add?
\end{itemize}

\noindent
Optionally, you may also download the source code, modify it, and run your own experiments.

\noindent\textbf{Good luck!}

\end{document}